\begin{abstract}
Agent-based program repair offers to
automatically resolve complex bugs end-to-end by
combining the planning, tool use, and code generation
abilities of modern LLMs. Recent work has explored
the use of agent-based repair approaches on 
the popular open-source SWE-Bench~\cite{swebench}, a collection of bugs from highly-rated GitHub Python projects.
In addition, various agentic approaches such as SWE-Agent~\cite{sweagent} have been proposed to solve bugs in this benchmark.

This paper explores the viability of using an agentic approach to address bugs in an enterprise context. 
To investigate this, we curate
an evaluation set of 178 bugs drawn from
Google's issue tracking system. This dataset
spans both human-reported (78) and machine-reported bugs (100).

To establish a repair performance baseline on this benchmark, we implement Passerine, an agent similar in spirit to SWE-Agent that can work within Google's development 
environment.  We show that with 20 trajectory samples and Gemini 1.5 Pro,
Passerine can produce a patch that passes bug tests (i.e., plausible) for 
73\% of machine-reported and 25.6\% of human-reported bugs 
in our evaluation set. 
After manual examination, we found that 43\% of machine-reported bugs and 17.9\% of human-reported bugs have at least one patch that is semantically equivalent to the ground-truth patch.

These results establish a baseline on an industrially
relevant benchmark, which as we show,
contains bugs drawn from
a different distribution---in terms of language diversity, size, and spread of changes, etc.---compared to those
in the popular SWE-Bench dataset.









\end{abstract}
